\chapter{Results}
\label{chap:results}

\section{Primary Accuracy Results}

\Cref{tab:primary-accuracy} reports mean accuracy and standard deviation from the primary
three-seed sweep.

\begin{table}[H]
\centering
\caption{Primary SST-2 accuracy (\%) over three seeds. Best PEFT result per budget is bold.}
\label{tab:primary-accuracy}
\small
\begin{tabular}{@{}rrrrrr@{}}
\toprule
Samples & Full & LoRA & AdaLoRA & Star-LoRA & Star Rank \\
\midrule
100   & 53.48$\pm$3.95 & \textbf{52.75$\pm$0.53} & 51.07$\pm$0.37 & 51.03$\pm$0.41 & 4.67 \\
500   & 73.66$\pm$0.58 & \textbf{67.47$\pm$5.86} & 57.19$\pm$5.40 & 58.37$\pm$4.58 & 6.33 \\
1,000 & 75.23$\pm$3.10 & \textbf{85.17$\pm$1.86} & 65.14$\pm$3.33 & 69.34$\pm$1.97 & 7.00 \\
5,000 & 77.45$\pm$1.73 & 91.25$\pm$0.33 & 86.58$\pm$0.57 & \textbf{91.55$\pm$0.07} & 9.00 \\
10,000& 76.83$\pm$0.50 & \textbf{92.28$\pm$0.65} & 90.02$\pm$0.60 & 91.97$\pm$0.64 & 10.00 \\
\bottomrule
\end{tabular}
\end{table}

Star-LoRA's clearest advantage is relative to AdaLoRA.
\begin{table}[H]
\centering
\caption{Star-LoRA accuracy difference relative to standard AdaLoRA.}
\label{tab:adalora-delta}
\begin{tabular}{@{}rrl@{}}
\toprule
Samples & Difference (percentage points) & Interpretation \\
\midrule
100 & $-0.04$ & no meaningful improvement \\
500 & $+1.18$ & small improvement \\
1,000 & $+4.20$ & substantial improvement \\
5,000 & $+4.97$ & largest improvement \\
10,000 & $+1.95$ & moderate improvement \\
\bottomrule
\end{tabular}
\end{table}

At 100 examples, both adaptive methods remain near chance-level accuracy and have low macro F1.
The run contains only 39 estimated steps, leaving little evidence for dynamic allocation. At 500
examples, Star-LoRA begins to improve AdaLoRA but remains far behind fixed LoRA and full
fine-tuning. At 1,000 examples the gap over AdaLoRA grows, although LoRA is still much stronger.
At 5,000 examples Star-LoRA slightly exceeds LoRA and clearly exceeds AdaLoRA. At 10,000 examples
LoRA regains a small lead.

\section{Macro F1}

\begin{table}[H]
\centering
\caption{Primary SST-2 macro F1 (\%) over three seeds.}
\label{tab:primary-f1}
\small
\begin{tabular}{@{}rrrrr@{}}
\toprule
Samples & Full & LoRA & AdaLoRA & Star-LoRA \\
\midrule
100   & 42.71$\pm$13.33 & \textbf{49.03$\pm$3.78} & 39.58$\pm$7.18 & 39.56$\pm$7.20 \\
500   & \textbf{73.54$\pm$0.49} & 66.73$\pm$5.97 & 56.11$\pm$6.62 & 57.65$\pm$5.18 \\
1,000 & 75.02$\pm$3.40 & \textbf{85.15$\pm$1.87} & 64.98$\pm$3.22 & 69.16$\pm$2.11 \\
5,000 & 77.32$\pm$1.88 & 91.24$\pm$0.34 & 86.57$\pm$0.57 & \textbf{91.55$\pm$0.06} \\
10,000& 76.76$\pm$0.45 & \textbf{92.27$\pm$0.65} & 90.02$\pm$0.60 & 91.97$\pm$0.64 \\
\bottomrule
\end{tabular}
\end{table}

Accuracy and macro F1 are nearly identical once performance becomes high, which is expected for a
roughly balanced binary task. At 100 examples the difference is larger, indicating unstable or
class-skewed predictions. Extreme low-resource runs are not favorable to the current adaptive
mechanism.

\section{Rank Selection}

\begin{figure}[H]
  \centering
  \fbox{\begin{minipage}{0.78\textwidth}
  \centering
  \begin{tabular}{c|ccccc}
  Training examples & 100 & 500 & 1,000 & 5,000 & 10,000\\
  \hline
  Mean selected rank & 4.67 & 6.33 & 7.00 & 9.00 & 10.00
  \end{tabular}\\[1em]
  $\text{Data budget}\ \nearrow \quad\Longrightarrow\quad
  \text{selected adapter capacity}\ \nearrow$
  \end{minipage}}
  \caption{The deterministic policy increases target rank with the sample budget.}
  \label{fig:rank-trend}
\end{figure}

The monotonic trend demonstrates that the policy operates as designed. It does not prove that each
rank is optimal. Establishing optimality requires a controlled rank sweep under identical modules
and schedules.

\section{Stability Across Seeds}

At 5,000 examples, Star-LoRA records an accuracy standard deviation of $0.0007$, compared with
$0.0057$ for AdaLoRA and $0.0033$ for LoRA. In percentage-point terms these are 0.07, 0.57, and
0.33. Star-LoRA combines its largest AdaLoRA improvement with its smallest observed variance.

This result should not be generalized to every budget. At 100 examples Star-LoRA's macro-F1
standard deviation is 7.20 percentage points, and at 500 examples accuracy varies by 4.58 points.
The method appears stable in the 5,000-example condition, not uniformly stable.

\section{Runtime and Throughput}

\begin{table}[H]
\centering
\caption{Primary mean training runtime in seconds.}
\label{tab:runtime}
\begin{tabular}{@{}rrrrr@{}}
\toprule
Samples & Full & LoRA & AdaLoRA & Star-LoRA \\
\midrule
100 & 8.44 & 8.65 & 11.77 & 15.72 \\
500 & 19.38 & 17.89 & 25.97 & 56.33 \\
1,000 & 32.97 & 29.42 & 43.10 & 100.78 \\
5,000 & 142.92 & 120.93 & 183.26 & 444.96 \\
10,000 & 277.48 & 237.05 & 362.96 & 851.83 \\
\bottomrule
\end{tabular}
\end{table}

At 5,000 examples Star-LoRA is approximately $2.43\times$ slower than AdaLoRA and
$3.68\times$ slower than LoRA. The 0.30 percentage-point advantage over LoRA at that budget is
therefore obtained at a large runtime premium. At 10,000 examples Star-LoRA is slower and slightly
less accurate than LoRA.

The primary reason is visible in the run metadata. Standard AdaLoRA targets
\texttt{q\_proj} and \texttt{v\_proj}. The 5,000-example Star-LoRA run selects rank 9 and targets
all four attention projections plus \texttt{gate\_proj}, \texttt{up\_proj}, and
\texttt{down\_proj}. The method increases both rank and adapter surface.

\section{Secondary Five-Seed Sweep}

\begin{table}[H]
\centering
\caption{Accuracy (\%) in the separate five-seed archive. This sweep is not pooled with the
primary experiment.}
\label{tab:secondary}
\begin{tabular}{@{}rrrrr@{}}
\toprule
Samples & Full & LoRA & AdaLoRA & Star-LoRA \\
\midrule
100 & 81.88 & 50.02 & 50.21 & 50.21 \\
500 & 86.79 & 53.23 & 50.41 & 50.50 \\
1,000 & 88.14 & 63.81 & 52.94 & 53.83 \\
5,000 & 89.47 & 90.02 & 71.33 & 84.01 \\
10,000 & 88.44 & 91.61 & 79.86 & 89.84 \\
\bottomrule
\end{tabular}
\end{table}

This sweep presents a different absolute performance regime, especially for full fine-tuning,
which confirms that archived experiment groups should not be mixed without complete configuration
records. Within the sweep, Star-LoRA again improves AdaLoRA most strongly at larger budgets: by
12.68 percentage points at 5,000 examples and 9.98 points at 10,000. LoRA remains the best PEFT
method at both budgets.

\section{Layer-Wise Heuristic}

\begin{table}[H]
\centering
\caption{Five-seed global and layer-wise heuristic results from separate archived runs.}
\label{tab:layerwise}
\begin{tabular}{@{}rrrrr@{}}
\toprule
Samples & Global Acc. & Layer Acc. & Global Time & Layer Time \\
\midrule
100 & 50.05 & 50.07 & 14.17 & 11.96 \\
500 & 49.63 & 49.70 & 43.15 & 29.91 \\
1,000 & 51.01 & 51.08 & 71.69 & 49.25 \\
5,000 & 71.93 & 68.28 & 310.44 & 209.98 \\
\bottomrule
\end{tabular}
\end{table}

At 5,000 examples the layer-wise heuristic reduces runtime by about 32\% but lowers accuracy by
3.65 percentage points. Concentrating additional adapters in upper layers is an efficiency
trade-off, but the chosen thresholds do not preserve the global policy's accuracy.

\section{Exploratory LLM-Assisted Allocation}

\begin{table}[H]
\centering
\caption{Five-seed exploratory results at 5,000 examples.}
\label{tab:llm}
\begin{tabular}{@{}lrrr@{}}
\toprule
Method & Accuracy (\%) & Macro F1 (\%) & Runtime (s) \\
\midrule
Global rule-based & 71.93 & 71.88 & 310.44 \\
Layer-wise heuristic & 68.28 & 68.21 & 209.98 \\
Claude layer-wise & 67.45 & 67.40 & 190.99 \\
\bottomrule
\end{tabular}
\end{table}

The Claude policy is about 38.5\% faster than the global rule-based method but 4.48 percentage
points less accurate. It is also 0.83 points less accurate than the deterministic layer-wise
heuristic while being about 9\% faster. The generated layout selected a cheaper adapter surface,
but did not identify the modules needed to preserve task quality.

No Gemini accuracy is reported. The ``Gemini'' and heuristic rows in the pilot summary are
identical because the Gemini API calls returned quota errors and the implementation fell back to
the heuristic. Treating these rows as independent methods would double-count one policy.

\section{Answers to the Research Questions}

\paragraph{RQ1.}
Star-LoRA improves standard AdaLoRA at four of five budgets in the primary sweep, with the largest
gain of 4.97 percentage points at 5,000 examples.

\paragraph{RQ2.}
The selected rank increases from approximately 5 to 10 as the training budget grows from 100 to
10,000 examples. The policy is responsive, although optimality has not been established.

\paragraph{RQ3.}
Star-LoRA is competitive with LoRA at 5,000 and 10,000 examples, but LoRA is substantially better
at 500 and 1,000 examples and slightly better at 10,000.

\paragraph{RQ4.}
Dataset-aware module expansion is expensive. At 5,000 examples the primary Star-LoRA run is
approximately 2.43 times slower than AdaLoRA. Seed variance is low in this condition but not in
extreme low-resource conditions.

\paragraph{RQ5.}
The Claude-generated policy reduces runtime but does not improve accuracy. Gemini attempts are
invalid for performance comparison because they used heuristic fallback.
